% ============================================================
%  CHƯƠNG 1: TỔNG QUAN VỀ LĨNH VỰC NGHIÊN CỨU
%  Heading 1: Times New Roman 14pt, đậm, in hoa
%  Được gọi từ main.tex với \chapter{TỔNG QUAN VỀ LĨNH VỰC NGHIÊN CỨU}
% ============================================================

\section{Giới thiệu tổng quan}
% Heading 2: Times New Roman 13pt, đậm

[Giới thiệu tổng quan về lĩnh vực nghiên cứu. Nêu tầm quan trọng của chủ đề
và xu hướng phát triển hiện nay.]

% -------------------------------------------
\subsection{Tiểu mục thứ nhất}
% Heading 3: Times New Roman 13pt, đậm nghiêng

[Nội dung tiểu mục. Văn bản nội dung sử dụng Times New Roman 13pt, giãn dòng 1.5,
căn lề hai bên, thụt đầu dòng 1cm.]

\subsubsection{Tiểu mục thứ hai}
% Heading 4: Times New Roman 13pt, nghiêng

[Nội dung tiểu mục thứ hai. Đây là cấp sâu nhất được sử dụng (tối đa 4 chữ số: 1.1.1.1).]

% -------------------------------------------
\section{Tổng quan các công trình nghiên cứu}

\subsection{Các nghiên cứu trong nước}

[Tổng hợp và phân tích các công trình trong nước có liên quan. Mỗi công trình cần:
tên tác giả, năm công bố, phương pháp sử dụng, kết quả đạt được, hạn chế.]

\subsection{Các nghiên cứu ngoài nước}

[Tổng hợp các công trình quốc tế tiêu biểu, ưu tiên các bài báo trong 5 năm gần nhất
được đăng trên các tạp chí/hội nghị uy tín (Scopus, ISI).
Có thể trích dẫn theo chuẩn IEEE: \cite{example_ref1}, \cite{example_ref2}.]

% Ví dụ bảng tổng hợp nghiên cứu
\begin{table}[h!]
\centering
\fontsize{13pt}{14.95pt}\selectfont
\begin{tabularx}{\textwidth}{|c|>{\raggedright\arraybackslash}X|
                               >{\raggedright\arraybackslash}X|c|}
\hline
\textbf{STT} & \textbf{Nghiên cứu} & \textbf{Phương pháp chính} & \textbf{Năm} \\
\hline
1 & Tác giả A và cộng sự -- Mô tả ngắn nội dung nghiên cứu
  & RandomForest, XGBoost & 2022 \\
\hline
2 & Tác giả B và cộng sự -- Mô tả ngắn nội dung nghiên cứu
  & LSTM, Transformer     & 2023 \\
\hline
3 & Tác giả C và cộng sự -- Mô tả ngắn nội dung nghiên cứu
  & BERT, GPT             & 2024 \\
\hline
\end{tabularx}
\caption{Tóm tắt các nghiên cứu liên quan}
\label{tab:tong_quan_nghien_cuu}
\end{table}

\subsection{Nhận xét và khoảng trống nghiên cứu}

[Đánh giá ưu điểm và hạn chế của các nghiên cứu trên.
Xác định rõ \textbf{research gap} mà luận văn này sẽ giải quyết.]

% -------------------------------------------
\section{Ví dụ chèn hình ảnh}

[Hình \ref{fig:quy_trinh_nghien_cuu} minh họa quy trình nghiên cứu tổng quát
của luận văn.]

\begin{figure}[H]
  \centering
  \includegraphics[width=0.8\textwidth]{placeholder_figure}  % thay bằng file ảnh thực tế
  \caption{Quy trình nghiên cứu tổng quát}
  \label{fig:quy_trinh_nghien_cuu}
\end{figure}

% -------------------------------------------
\section{Ví dụ công thức toán học}

[Công thức \eqref{eq:rmse} trình bày độ đo RMSE được sử dụng để đánh giá mô hình:]

\begin{equation}
  \text{RMSE} = \sqrt{\frac{1}{n} \sum_{i=1}^{n} \left( \hat{y}_i - y_i \right)^2}
  \label{eq:rmse}
\end{equation}

Trong đó: $\hat{y}_i$ là giá trị dự đoán; $y_i$ là giá trị thực tế;
$n$ là số lượng mẫu trong tập kiểm tra.
